\section{Conclusion and Limitations}

This paper studied online local planning for mobile-base open-liquid transport and proposed \SPMPC{}, a slosh-aware planning method that embeds low-order slosh modal states into an \MPCC{} formulation.
The central premise is that liquid sloshing is not only a trajectory-smoothing issue, but a state-prediction problem with dynamic memory.
If a planner penalizes only velocity, acceleration, or command variation, it may fail to represent the effect of the current liquid phase on future response.
\SPMPC{} forms an augmented state from base motion, path progress, and low-order slosh variables, jointly optimizes contouring error, lag error, path progress, control magnitude, control smoothness, and predicted slosh response, and executes the first optimized base command in receding horizon.
The paper also formulates a soft Slosh-risk Reference Governor as a pre-\MPCC{} finite-horizon reference-adaptation module that contracts the speed reference without directly rewriting the final base command.

Representative fixed-path physical trials provide initial evidence for the main claim.
For \Bzero{}, \Bsmooth{}, and \Bslosh{}, three trials per method reached the goal.
Over the 10--90\% path-progress window, \Bslosh{} reduced the model-proxy p95 and RMS by 36.4\% and 37.4\% relative to the smooth-only baseline, and reduced RGB $H_{\mathrm{vis}}$ p95 and RMS by 29.2\% and 15.9\%.
These p95/RMS trends indicate that the slosh-aware variant can provide liquid-response improvement beyond smooth-only control in the representative trials.
The evidence is deliberately interpreted conservatively: RGB RMS decreases less than the model proxy, \Bslosh{} takes longer to reach the goal, and its path-projection and angular-velocity metrics are not uniformly best.
Internal model proxies are therefore treated as control-side diagnostic quantities rather than ground-truth liquid height, complete candidate-ranking criteria, or hard safety signals.

The conclusions hold under the assumptions of a precomputed geometrically feasible reference path, low-order slosh modeling, and external liquid-response evaluation.
\SPMPC{} is positioned as a method for moving liquid dynamic memory into the online \MPCC{} local-planning layer of a standard mobile base, not as a high-fidelity free-surface reconstruction method or a formal spill-free controller.
The Slosh-risk Reference Governor is likewise a soft reference-adaptation mechanism, not a hard safety interlock or a spill-prevention proof.
Future extensions include broader validation across paths, fill levels, container geometries, ordinary local-planner baselines, mobile-base anti-slosh references, and governor ablations, as well as measured container acceleration, real liquid-surface feedback, and uncertainty-aware constraints for more complex navigation scenarios.
